\documentclass[11pt,a4paper,fontset=fandol]{ctexart}
\usepackage[margin=20mm]{geometry}
\usepackage{amsmath,amssymb,booktabs,array,tabularx,enumitem,setspace,xcolor,hyperref,fancyhdr}
\usepackage{xurl}
\hypersetup{colorlinks=true,linkcolor=blue!45!black,urlcolor=blue!45!black,citecolor=blue!45!black,pdftitle={Task II Step 3：真实模型试算与共享方法回查}}
\setstretch{1.10}
\setlength{\emergencystretch}{3em}
\setlength{\headheight}{15pt}
\setlength{\tabcolsep}{5pt}
\renewcommand{\arraystretch}{1.17}
\setlist{nosep,leftmargin=2em}
\pagestyle{fancy}\fancyhf{}
\fancyhead[L]{\small Task II：真实模型试算}
\fancyhead[R]{\small Step 3，完成待审阅}
\fancyfoot[C]{\thepage}
\newcommand{\qs}{Q_{\mathrm S}}
\newcommand{\qr}{Q_{\mathrm R}}
\newcommand{\RI}{\mathrm{RI}}
\newcommand{\op}{\mathrm{op}}
\newcommand{\loc}{\mathrm{ports}}
\newcommand{\cn}[1]{\left\lceil #1/128\right\rceil}
\newcommand{\para}[1]{\par\smallskip\noindent\textbf{#1}\quad}
\newcommand{\Byte}{\,\mathrm{Byte}}
\newcommand{\KiB}{\,\mathrm{KiB}}
\newcommand{\MiB}{\,\mathrm{MiB}}
\title{\bfseries 真实模型试算与共享方法回查\\[5pt]\large 两个完整实例与一个不等宽 head 定点检查}
\author{}\date{2026 年 9 月 24 日}
\begin{document}
\maketitle\vspace{-1.6em}
\noindent\textbf{结论：本轮 22 个汇总工况支持保持 Step 1 提取与 Step 2 方法。}
Ministral 3 8B（2512）和 Qwen3.6-35B-A3B 各完成 10 个工况；MiMo-V2.5-Pro 仅完成 $L=1024$ 的两个 Attention 工况。
以本地 HEAD / Step 2 审阅提交 \texttt{2eb7e7c} 为起点，无权重下载、模型运行或性能测量。

\section{先看实际结果：比例相同也有信息}
以下均为精确 RI。FFN 顺序为 $B=(8,64,512)$；Attention 顺序为
$L=(1024,16384,131072)$，即 $(1\mathrm K,16\mathrm K,128\mathrm K)$，$\mathrm K=1024$。
QKV 为一个输入 token，权重在窗口前驻留。

\para{主参考：局部接收端累计 ports}
\begin{center}\small\begin{tabular}{@{}lll@{}}\toprule
Workload & Ministral 3 8B (2512) & Qwen3.6-35B-A3B\\\midrule
QKV Projection & $\infty$ & $\infty$\\[4pt]
FFN / MoE & $\frac{1}{16},\;\frac{1}{2},\;4$ & $\frac{1}{16},\;\frac{1}{2},\;4$\\[4pt]
Attention Prefill & $\frac{2177}{128},\;\frac{32897}{128},\;\frac{262273}{128}$ & $\frac{2177}{64},\;\frac{32897}{64},\;\frac{262273}{64}$\\[4pt]
Attention Decode & $32,\;512,\;4096$ & $64,\;1024,\;8192$\\[4pt]
\bottomrule\end{tabular}
\end{center}
\para{对照：整矩阵阶段入口 operator}
\begin{center}\small\begin{tabular}{@{}lll@{}}\toprule
Workload & Ministral 3 8B (2512) & Qwen3.6-35B-A3B\\\midrule
QKV Projection & $\infty$ & $\infty$\\[4pt]
FFN / MoE & $\frac{3}{3584},\;\frac{3}{448},\;\frac{3}{56}$ & $\frac{5}{768},\;\frac{5}{96},\;\frac{5}{12}$\\[4pt]
Attention Prefill & $\frac{1281}{128},\;\frac{16641}{128},\;\frac{131329}{128}$ & $\frac{1537}{128},\;\frac{16897}{128},\;\frac{131585}{128}$\\[4pt]
Attention Decode & $18,\;258,\;2050$ & $20,\;260,\;2052$\\[4pt]
\bottomrule\end{tabular}
\end{center}

两款 FFN 的局部 RI 均为 $B/128$，但 Dense 三矩阵容量为 168 MiB，单专家为 3 MiB，相差 56 倍。
同一 $L$ 下，两款 Attention 的 ports 输入与调用数完全相同，Qwen 的 KV 写入只有一半，因而局部 RI 为两倍。
这些抵消来自真实维度与所声明映射，不是遗漏模型差异。

\para{阅读与数据边界}
本稿给出四类代表工况的中文代入及 MiMo 尾块检查；全部分项的两套
$(\qs,\qr,\RI)$、调用、容量与写入片保存在 \path{data/results.json} 和 \path{data/results.csv}。
Byte 与分数为原始精确量，显示单位单独转换，$1\KiB=1024\Byte$、$1\MiB=2^{20}\Byte$。
各行窗口不同，不能直接相加作为端到端需求。

\clearpage
\section{原始结构回查与计算接口}
\para{固定模型、层与实现}
层号从 0 起。模型 revision 如下，未换 checkpoint：
\begin{quote}\small
Ministral 3 8B Base 2512，第 0 层 full GQA + Dense：\\
\texttt{d4883f9b36aa2e5d775730d3fdba3d30de51a8ef}\\
Qwen3.6-35B-A3B，第 3 层 full GQA + MoE：\\
\texttt{995ad96eacd98c81ed38be0c5b274b04031597b0}\\
MiMo-V2.5-Pro，第 7 层 global GQA：\\
\texttt{21d1ecfecd7bd70f31be25ca49d7edd21f003659}
\end{quote}
前两者的 Transformers 静态实现固定于
\texttt{2c4914fb939fe9de0d8e7a798af4684d552f18b4}；MiMo remote code 与 checkpoint 同 revision。
共享方法沿用 \texttt{WS128-INT8-semantic-banks-v1}，未修改版本。

\begin{center}\small
\begin{tabular}{@{}lrrrrrr@{}}\toprule
模型 & $D$ & $H_q$ & $H_{kv}$ & $d_{QK}$ & $d_V$ & $F$\\\midrule
Ministral & 4096 & 32 & 8 & 128 & 128 & 14336\\
Qwen3.6 & 2048 & 16 & 2 & 256 & 256 & 512\\
MiMo（仅 Attention） & 6144 & 128 & 8 & 192 & 128 & 不求值\\\bottomrule
\end{tabular}\end{center}

\para{不是只相信 models.json}
Ministral 的 HF \texttt{text\_config} 与原生 \texttt{params.json} 逐键一致；原生
\texttt{v\_head\_dim=null} 在固定实现中对应 V 采用 128，未按零宽处理。
Qwen 配置第 3 层是 \texttt{full\_attention}，原生 \texttt{q\_proj} 输出
$16\times256\times2=8192$，按每头切开 query 与 gate；专家张量为
$[256,1024,2048]$ 的 gate/up 和 $[256,2048,512]$ 的 down。
主对象只取其中一个专家，不乘 8 个路由命中或 256 个专家。
MiMo 的 \texttt{hybrid\_layer\_pattern[7]=0}，全局层没有 sink；192 是完整 QK 维度，不能以 64 维部分 RoPE 或 $D/H_q=48$ 替代。
证据定位见文末及 \path{data/checks.json}。

\para{角色与交接}
K/V 在 GQA 扩展前写入，每 KV head 一份：Ministral 保存 RoPE 后 K；Qwen 保存归一化、RoPE 后 K；MiMo 保存部分 RoPE 后的完整 192 维 K 与已乘 0.612 的 V。
每个 head 的 K 排为 $i\times d_{QK}$（token 沿输出行），V 排为 $d_V\times i$（token 沿输入列）。
软件 cache 布局不直接规定 CIM 的矩阵方向。

各角色参考均为 INT8 / 1 Byte；原生 BF16/FP8 记录未覆盖，也不作精度验证。
语义矩阵独立分块，源代码融合仅需切片或重排，不合并接收端、不复制权重。
部分和、norm/RoPE、softmax、SiLU 与逐元素门控在数字接口完成；只有成为下一次矩阵输入的 down 向量或 Attention 系数再计入 $\qs$。
所有计算假定状态有足够驻留空间，未产生容量重载。

\clearpage
\section{QKV：真实矩形、额外 gate 与静态端点}
\para{Ministral：原生结构到分块}
输入为 $x\in\mathbb Z^{4096}$，resident 权重分别为
$W_Q:4096\times4096$，$W_K,W_V:1024\times4096$。
每个输入维分为 32 个 128 维片；输出块数为 $32,8,8$，所有尾块均完整。
一个 token 使每个 tile 求值一次。
\begin{center}\small
\begin{tabular}{@{}lrrrr@{}}\toprule
分项 & tile 数/调用数 & $\qs^{\loc}$ (Byte) & 独立入口 $\qs^{\op}$ (Byte) & $\qr$ (Byte)\\\midrule
Q & 1024 & 131072 & 4096 & 0\\
K & 256 & 32768 & 4096 & 0\\
V & 256 & 32768 & 4096 & 0\\\bottomrule
\end{tabular}\end{center}
于是两套汇总分别为
\begin{align*}
(\qs^{\loc},\qr^{\loc},\RI^{\loc})
 &=\bigl[4096(32+8+8),\;0,\;\infty\bigr]
 =(196608,0,\infty),\\
(\qs^{\op},\qr^{\op},\RI^{\op})
 &=(3\times4096-2\times4096,\;0,\;\infty)=(4096,0,\infty).
\end{align*}
三分项 operator 输入是同一 $x$，汇总扣除 8192 Byte 重叠；ports 的不同接收端直接相加。
权重有效容量与 tile 分配容量均为 $1536\times16384=25165824$ Byte（24 MiB），但本窗口写入为零。

\para{Qwen：完整保留 output gate}
语义矩阵 Q、G、K、V 分别为
$4096\times2048,4096\times2048,512\times2048,512\times2048$。
Q/G 在原生 q-projection 中按 head 打包；独立检查验证重排后全部行恰好覆盖一次。
输入分 16 块，输出块数为 $32,32,4,4$：
\begin{center}\small
\begin{tabular}{@{}lrrrr@{}}\toprule
分项 & tile 数/调用数 & $\qs^{\loc}$ (Byte) & 独立入口 $\qs^{\op}$ (Byte) & $\qr$ (Byte)\\\midrule
Q & 512 & 65536 & 2048 & 0\\
G（output gate） & 512 & 65536 & 2048 & 0\\
K & 64 & 8192 & 2048 & 0\\
V & 64 & 8192 & 2048 & 0\\\bottomrule
\end{tabular}\end{center}
ports 汇总为 $(147456,0,\infty)$；operator 从 $4\times2048$ 扣除
$3\times2048$，为 $(2048,0,\infty)$；共 1152 次调用，权重容量 18 MiB。
若错误删去 G，会少 65536 Byte 局部输入、512 次调用与 8 MiB resident 权重，静态 RI 却仍是 $\infty$。
因此只检查 RI 无法发现这一结构遗漏。

\para{为什么需求变而比例不变}
两款 QKV 的输入、权重规模、调用数都不同；$\qr=0$ 仅来自权重预驻留的窗口。
KV 建立属于 Attention 行，不是这里的权重更新；G 是后续输出门控，不是 query 或 KV head。
本行不含输出投影 $W_O$，也不拆成 Prefill/Decode 两行。

\clearpage
\section{FFN / 单路由专家：一次装载、服务 B 个向量}
\para{Ministral，$B=8$ 的完整代入}
原生 SwiGLU 的 $W_g,W_u$ 为 $14336\times4096$，$W_d$ 为
$4096\times14336$。前两者各 $112\times32$ 个 tile，down 为 $32\times112$。
三者权重在本窗口各装载一次，随后留驻服务 8 个向量。
$z=\operatorname{SiLU}(W_gx)\odot(W_ux)$ 重定标后是 down 的 streaming 输入。
\begin{center}\small
\begin{tabular}{@{}lrrrr@{}}\toprule
矩阵 & $\qs^{\loc}$ (Byte) & $\qs^{\op}$ (Byte) & $\qr$ (Byte) & 调用数\\\midrule
gate & $8\times112\times4096=3670016$ & 32768 & 58720256 & 28672\\
up & 3670016 & 32768 & 58720256 & 28672\\
down & $8\times32\times14336=3670016$ & 114688 & 58720256 & 28672\\\bottomrule
\end{tabular}\end{center}
\begin{align*}
(\qs^{\loc},\qr^{\loc},\RI^{\loc})&=(11010048,176160768,1/16),\\
(\qs^{\op},\qr^{\op},\RI^{\op})&=(147456,176160768,3/3584).
\end{align*}
operator 汇总 $8(4096+14336)=147456$ Byte；gate/up 共享 $x$，扣除一份
$8\times4096$，但 down 的 114688 Byte 不能再去掉。
共 10752 个 resident tile、86016 次求值；有效/分配容量均为 168 MiB。

\para{Qwen，单专家 $B=8$}
打包 gate/up 切开为两个 $512\times2048$，down 为 $2048\times512$。
各矩阵 64 个 tile、$1048576$ Byte 权重，8 向量各产生
$65536$ Byte 局部输入和 512 次调用。总计
\begin{align*}
(\qs^{\loc},\qr^{\loc},\RI^{\loc})&=(196608,3145728,1/16),\\
(\qs^{\op},\qr^{\op},\RI^{\op})&=(8(2048+512),3145728,5/768).
\end{align*}
operator 分项为 $(16384,16384,4096)$ Byte，扣除共享输入 16384 Byte。
192 个 resident tile 占 3 MiB，求值 1536 次。路由器、共享专家及其输出门不属于这个单专家对象。

\para{三档 B：变化与不变}
局部输入与调用随 $B$ 成比例；一次权重装载、最终容量、tile 数均不变。
两模型 $D,F$ 都整除 128，因此每个矩阵均有
$\qs^{\loc}=BDF/128$、$\qr=DF$，汇总也为
\[
\RI^{\loc}=\frac{3BDF/128}{3DF}=\frac B{128}
=\left(\frac1{16},\frac12,4\right).
\]
Dense 与单专家的 ports 输入、写入、tile 数和调用数均相差 56 倍，这个因子在比值中抵消。
operator 则为 $B(D+F)/(3DF)$，没有相同抵消。
这里的相同局部 RI 依赖整除、三矩阵独立分块、等角色 Byte 和同一驻留期；不推广为所有 FFN 的无条件恒等式。

\clearpage
\section{Attention Prefill：从空状态建立因果前缀}
\para{Ministral，$L=1024$ 的完整代入}
8 个 KV head 各持一份 K 和 V，每份服务 4 个 query heads。
前缀 $i$ 的 resident K 为 $i\times128$，V 为 $128\times i$。
K 分为 $\lceil i/128\rceil\times1$ 个 tile；V 为 $1\times\lceil i/128\rceil$。
每个 token 先追加 K 的 $1\times128$ 行和 V 的 $128\times1$ 列，再求值全部 32 个头。
\[
S_{128}(1024)=128(1+2+\cdots+8)=4608,\qquad
P(1024)=\sum_{i=1}^{1024}i=524800.
\]
QK 的 query 在每个序列输出块重复接收，AV 的系数长度为 $i$：
\begin{align*}
Q_{\mathrm S,QK}^{\loc}&=32\times128\times4608=18874368,\\
Q_{\mathrm S,QK}^{\op}&=32\times128\times1024=4194304,\\
Q_{\mathrm S,AV}^{\loc}&=Q_{\mathrm S,AV}^{\op}=32\times524800=16793600,\\
Q_{\mathrm R,K}&=Q_{\mathrm R,V}=1024\times8\times128=1048576.
\end{align*}
因此，Byte 三元组和 RI 为
\begin{align*}
(\qs^{\loc},\qr^{\loc},\RI^{\loc})&=(35667968,2097152,2177/128),\\
(\qs^{\op},\qr^{\op},\RI^{\op})&=(20987904,2097152,1281/128).
\end{align*}
QK、AV 各调用 $32\times4608=147456$ 次，总计 294912。
最终各有 64 个 K/V tile，总容量 2 MiB；本窗口从空开始且只追加一次，所以累计写入刚好等于最终有效容量。
这不是把容量一般定义成 $\qr$。

\para{有效尾块不能被最终整除掩盖}
虽然最终 $L=1024$ 整除 128，前 127 个前缀以及每个块的增长期仍有尾块。
例如 $i=129$ 时 K 的输出块为 $(128,1)$，每块仍接收完整 128 维 query；V 的输入块为 $(128,1)$，第二次调用只接收 1 个有效系数。
这解释了 $S_{128}(L)$ 与 $P(L)$ 是不同累计量，不能用同一个近似代替。

\para{Qwen 与 Ministral 为什么输入相同}
Qwen 的 $H_qd_{QK}=16\times256=4096$，与 Ministral 相同；
$H_q\lceil d_V/128\rceil=16\times2=32$ 也相同。
故两者对任意同一 $L$ 的 ports QK/AV 输入逐项相等。
调用因子也相同：$H_q\lceil d_{QK}/128\rceil=32$ 和
$H_q\lceil d_V/128\rceil=32$。
但 Qwen 每 token 只写 $2(256+256)=1024$ Byte，Ministral 为 2048 Byte。
在 1K 下，Qwen 的三元组为 ports $(35667968,1048576,2177/64)$，
operator $(12591104,1048576,1537/128)$。

\para{大长度计算}
一般令 $L=128m+r$，用
$S_{128}(L)=128m(m+1)/2+(m+1)r$ 和 $P(L)=L(L+1)/2$。
128K 两模型的 ports 输入均为 550026346496 Byte，调用数 4299161600；
Ministral/Qwen 的 KV 写入分别 256/128 MiB。脚本使用整数与紧凑求和，不创建 $L\times L$ 矩阵。

\clearpage
\section{Attention Decode：只收费本次追加，并比较窗口}
\para{Ministral，可见 $L=1024$}
初始已有 1023 个 token；本次追加 K 的 8 份 $1\times128$ 和 V 的 8 份
$128\times1$，分别写 1024 Byte。
最终每头 K 为 $1024\times128$、V 为 $128\times1024$，各 8 个 tile。
32 个 query heads 的一次求值为
\begin{center}\small
\begin{tabular}{@{}lrrrr@{}}\toprule
分项 & $\qs^{\loc}$ (Byte) & $\qs^{\op}$ (Byte) & 本次 $\qr$ (Byte) & 调用数\\\midrule
QK & $32\times8\times128=32768$ & 4096 & 1024 & 256\\
AV & $32\times1024=32768$ & 32768 & 1024 & 256\\\bottomrule
\end{tabular}\end{center}
\begin{align*}
(\qs^{\loc},\qr^{\loc},\RI^{\loc})&=(65536,2048,32),\\
(\qs^{\op},\qr^{\op},\RI^{\op})&=(36864,2048,18).
\end{align*}
初始有效容量为 2095104 Byte，最终为 2097152 Byte，分配容量也为 2 MiB；
$\qr=2048$ Byte 不是整个历史 cache 的重写。

\para{Qwen 的不同方向与共享}
K 为 $1024\times256$，每 KV head $8\times2$ 个 tile；V 为 $256\times1024$，每头 $2\times8$。
16 个 query heads 仍在 QK/AV 各调用 256 次、各接收 32768 Byte。
追加 K 为两份 $1\times256$（各切两个 $1\times128$），V 为两份 $256\times1$（各切两个 $128\times1$）。
总写入 1024 Byte、最终容量 1 MiB：ports 三元组 $(65536,1024,64)$；
operator $(20480,1024,20)$。

\para{同一 L，Prefill 与 Decode 的定量关系}
Prefill 写入是 Decode 的 $L$ 倍；Prefill 输入是从 1 到 $L$ 每次 Decode 输入的和，
因此 Prefill RI 是这些单步 RI 的均值（这里每步追加字节相等），而不是最后一步的 RI。
对两款主模型的固定 $L$（均为 128 的倍数），$d_{QK}=d_V=d$ 且 $d$ 整除 128，令 $g=H_q/H_{kv}$：
\begin{align*}
\RI^{\loc}_{\rm pre}&=g\frac{2L+129}{512},&
\RI^{\loc}_{\rm dec}&=g\frac{L}{128},&
\frac{Q_{\mathrm S,pre}^{\loc}}{Q_{\mathrm S,dec}^{\loc}}&=\frac{2L+129}{4},\\
\RI^{\op}_{\rm pre}&=\frac g2+\frac{g(L+1)}{4d},&
\RI^{\op}_{\rm dec}&=\frac g2+\frac{gL}{2d}.&&
\end{align*}
1K 的 ports 比值为 $\RI_{\rm pre}/\RI_{\rm dec}=2177/4096$；
输入累计比为 $2177/4=544.25$，写入比为 1024。
故 Decode RI 约为 Prefill 的两倍，不能说它们总需求或执行时间相近。

\para{差量检查的适用范围}
逐项验证 $\mathcal D_{\rm pre}(L)-\mathcal D_{\rm pre}(L-1)=\mathcal D_{\rm dec}(L)$，
$\mathcal D$ 为两种 $\qs$、$\qr$ 或调用数，\textbf{不对 RI 相减}。
这一关系仅对应本参考的逐前缀调用与同一 append 组织；换成批量求值、重新分块、重载或复制状态时须重新计数，不能把差量恒等式当成所有实现的性能规律。

\clearpage
\section{MiMo 定点：192=128+64，保留有效字节和调用}
\para{原生 global GQA 与布局}
只检查第 7 层、$L=1024$。$H_q/H_{kv}=128/8=16$，
$d_{QK}=192,d_V=128$；原生融合 QKV 宽度为 $24576+1536+1024=27136$。
投影仅作结构核对，本轮不计算 MiMo QKV 或 FFN。
每 KV head 的 K 为 $1024\times192$，分成 $8\times2$ 个 tile，输入宽度为 $(128,64)$；
V 为 $128\times1024$，有 8 个 tile。
每次追加 K 切成 $1\times128$ 和 $1\times64$，V 为 $128\times1$。
8 个头合计 24 个逻辑写入片，不代表 24 次物理写事务。

\para{两个窗口的分项原始 Byte}
\begin{center}\small
\begin{tabular}{@{}llrrrr@{}}\toprule
窗口 & 分项 & $\qs^{\loc}$ & $\qs^{\op}$ & $\qr$ & 调用数\\\midrule
Prefill & QK & 113246208 & 25165824 & 1572864 & 1179648\\
 & AV & 67174400 & 67174400 & 1048576 & 589824\\\midrule
Decode & QK & 196608 & 24576 & 1536 & 2048\\
 & AV & 131072 & 131072 & 1024 & 1024\\\bottomrule
\end{tabular}\end{center}
Prefill 的 QK 输入为 $128\times(128+64)\times4608$，调用为
$128\times2\times4608$；AV 输入为 $128\times524800$，调用为 $128\times4608$。
Decode 的 QK 输入为 $128\times8\times(128+64)$，调用为 $128\times8\times2$；
AV 输入为 $128\times1024$，调用为 $128\times8$。
两种边界的汇总必须分别展示：
\begin{align*}
\text{Prefill ports: } &(180420608,2621440,2753/40),\\
\text{Prefill operator: } &(92340224,2621440,1409/40),\\
\text{Decode ports: } &(327680,2560,128),\\
\text{Decode operator: } &(155648,2560,304/5).
\end{align*}

\para{尾片没有消失，也没有补成有效 128 Byte}
Decode QK 的 2048 次调用中，1024 次接收 128 Byte，1024 次接收 64 Byte；
Prefill QK 的两类调用各为 589824 次。
若误将尾片补满，Decode 会多计 $1024\times64=65536$ Byte；
若错误省掉尾片调用，则会漏掉 1024 次服务。
64 Byte 只是有效逻辑输入长度，不足以推出完整硬件服务时间减半；Task I 吞吐保持原样。

最终有效 KV 为 $8\times1024\times(192+128)=2621440$ Byte（2.5 MiB）；
K 128 个 tile、V 64 个 tile，共分配 3145728 Byte（3 MiB）。
Prefill 累计写入 2.5 MiB，Decode 仅追加 2560 Byte；三个量不混用。

\para{维度与分块分别影响什么}
$H_q$ 放大 query/系数服务，$H_{kv}$ 决定唯一 KV 写入与容量。
$d_{QK}$ 决定 query 字节、K 宽度和输入分块；$d_V$ 决定 V 写入与 AV 的输出块重放。
MiMo 在固定整块 $L$ 下，QK 和 AV 的局部 Decode RI 均为 128，但两项需求不同且调用数为 2:1。
此处比例相同也不能抹去半宽输入片的资源影响。

\clearpage
\section{方法回查：哪些结论可继续使用}
\para{独立检查怎样避免自证}
主计算使用共享 \path{counting.py}；独立检查从原始配置读取维度，以三条路径建立预期值：
\begin{enumerate}
\item 对真实矩阵直接切行列区间，累加每块有效输入、resident 元素和调用；
以“源阶段、向量编号”去重 operator 的同源输入，保留 down 新阶段。
\item Attention 将循环反过来：一个绝对 token tile 从出现到 $L$ 的存活期中，
每种部分宽度出现一次，长满后在剩余前缀重复调用。按存活期累加 QK/AV tile，
不使用主程序的三角数或 ceiling 闭式生成预期值。
\item 三个模型的真实 head 维度下，在 $L=1,127,128,129$ 显式枚举
query-head / KV-head / tile / 输入坐标与追加坐标，确认共享状态和尾片；再检查前缀差量。
\end{enumerate}
22 个主工况的两边界与所有分项一致；12 个边界点一致；7 个主长度 Decode 的 Prefill 差量一致。
所有大长度仅使用标量循环或紧凑直方图，不输出逐元素 Attention 事件文件。
原始文件哈希、实现行号、独立结果及有效输入宽度对应调用数均保存在 \path{data/checks.json}。

\para{一个足够小的映射对照}
Qwen 1K Decode 主参考只保留 2 份 KV。若改成每个 query head 都持久保存一份 KV，
共有 16 份，输入与调用不变，但每步写入与容量均乘 $g=8$：
\begin{center}\small
\begin{tabular}{@{}lrrrr@{}}\toprule
映射 & ports $\qs$ (Byte) & ports $\qr$ (Byte) & ports RI & 有效 KV 容量\\\midrule
一份 KV / KV head（主参考） & 65536 & 1024 & 64 & 1 MiB\\
一份 KV / query head（对照） & 65536 & 8192 & 8 & 8 MiB\\\bottomrule
\end{tabular}\end{center}
operator 同样从 $(20480,1024,20)$ 变为 $(20480,8192,5/2)$。
这不是原计数错误；它换了状态复制方式，说明 GQA 共享带来的 RI 优势依赖该映射。
对照未替代主参考，也不以潜在并行性推断延迟收益。

\para{三类判断}
\textbf{明确计数错误：本轮未发现。} 固定配置支持 gate、专家宽度、原生 head 和层型，
独立切片支持共享公式、追加次数与尾块处理。
\textbf{合理参考映射的结果：} FFN 的 $B/128$、两主模型 Attention 同输入、逐前缀 RI 关系，均可作为后续统一需求计算的基线。
\textbf{换执行方式才改变的项：} 状态复制、容量重载、跨矩阵尾块合并、不同 Prefill 调度或实际写入粒度，必须声明后再求需求/服务代价。

\para{能支撑与尚不能支撑的后续结论}
“足够空间、独立矩阵、逐前缀”可描述本轮全部所选矩阵任务，未出现结构不可表达的情形。
它不证明所有硬件有足够容量，或都能有效执行 K 行 / V 列小片更新。
结果已按组件保留，可继续检查聚合资源、更新粒度与尾块服务；总 RI 不直接换算为单 macro 延迟、器件排名或性能分区。
本轮无需修改 Step 1/2、Table II(a) 或 Task I。Step 3 审阅通过后具备进入 Step 4 的方法与数据条件；本轮到此停止。

\clearpage
\section{证据定位与复算入口}
\para{本地原始证据}
下列路径相对任务根目录；原件只读，未执行模型实现。固定 revision、SHA-256 与抽查结果见
\path{table_IIb/pilot/data/config.json}、\path{data/checks.json}（后者相对 pilot）。
\begin{enumerate}\raggedright
\item Ministral：\path{literature/02_ministral3_8b/raw/config.json} 的
\texttt{text\_config}，及同目录 \path{params.json}；
\path{literature/00_shared/raw/transformers/modeling_ministral3.py}：
111--173 行 Attention 构造、RoPE 与 cache update；176--188 行三矩阵 MLP；
90--101 行 GQA 扩展与两次矩阵乘。
\item Qwen：\path{literature/03_qwen36_35b_a3b/raw/config.json} 的
head、\texttt{layer\_types}、\texttt{attn\_output\_gate}、专家键；
\path{literature/00_shared/raw/transformers/modeling_qwen3_5_moe.py}：
751--825 行 q-projection 双宽输出、逐头切分、norm/RoPE、cache、sigmoid gate；
845--882 行单专家 gate/up 打包与 down；904--923 行独立共享专家分支；
947 行起选中层型。
\item MiMo：\path{literature/06_mimo_v25_pro/raw/config.json} 的全局层型、192/128 head、
\texttt{attention\_value\_scale} 和 sink 配置；同目录
\path{modeling_mimo_v2.py}：215--288 行维度、融合 QKV 与 sink 条件，
304--317 行 Value 缩放、部分 RoPE 与 cache update，389--406 行 global/SWA 选择。
\end{enumerate}

\para{产物与生成关系}
\begin{center}\small
\begin{tabularx}{\textwidth}{@{}lX@{}}\toprule
pilot 内路径 & 内容\\\midrule
\path{data/config.json} & 模型版本、层、角色精度、共享方法、输入快照与哈希\\
\path{data/results.json} / \path{results.csv} & 22 个工况；每分项两套需求、调用、容量、分块/追加\\
\path{data/checks.json} & 原始结构审计、独立结果、边界事件和调用宽度分布\\
\path{PREVIEW.md} & 两款完整实例的四行结果、固定 B/L 顺序\\
\path{scripts/generate.py} & 共享主计算、精确数据和两张预览表的生成/一致性检查\\
\path{scripts/check.py} & 原始配置、真实 tile、存活期和小规模事件的独立检查\\
\path{scripts/build.sh} / \path{render.py} & 复算、独立 PDF 编译、页面渲染\\
\path{data/pdf_qa.json} & 最终 PDF 哈希、页数和逐页目视核验记录\\\bottomrule
\end{tabularx}\end{center}
上表 \path{results.csv} 实际位于 \path{data/}，\path{render.py} 位于 \path{scripts/}。
\path{tex/pilot.zh.tex} 是独立中文入口；\path{output/pilot.zh.pdf} 是最终研究稿。
分项 operator 是独立入口，聚合去重另存；不能直接相加得到 QKV/FFN 总量。

\para{从仓库根复算与编译}
\begin{quote}\small\ttfamily
TASK2\_PYTHON=/opt/anaconda3/bin/python \\
sh tasks/task2\_table\_II\_workloads/table\_IIb/pilot/scripts/build.sh
\end{quote}
执行时将环境变量与 \texttt{sh} 放同一命令行，或先导出该变量。
默认检查不覆盖 JSON/CSV/表格，显式 \texttt{--emit} 才刷新派生文件。
共享 28 项合成检查、Table II(a) 数据/表格一致性与 Step 1 原始来源检查也纳入构建。
渲染使用 \path{scripts/render.py}（pypdfium2），逐页图保存在局部忽略的 \path{tmp/pdfs/}。

\para{状态与保留}
\path{STEP3_REVIEW.zh.md} 位于任务根目录；README 和研究配置记录
“Step 2 已通过；Step 3 完成待审阅”。Step 1/2 的旧审阅原稿、共享文件中的当时状态作为历史保留，
不为状态更新重编历史 PDF。没有 Git 提交、推送或 Step 4 分工。
\end{document}
